\section{Additional computational results}\label{app:computation}

This appendix gives the details and tables that Section~\ref{sec:computation}
summarizes.

\paragraph{Implementation (Section~\ref{sec:comp-impl}).}
The implementation differs from the analysis in the following ways; none of
them affects validity.
\begin{enumerate}[label=(\alph*)]
\item All coordinates share the mesh $h_j=s2^{-j}$, and the analysis that
applies is Lemma~\ref{lem:commonmesh} with the Euclidean condition number
$\kappa$, not Theorem~\ref{thm:approx}.
\item The cap of trial $\mu$ is $100\cdot2^\mu\lceil\log_2(n+2)\rceil$,
$12.5$ times the cap of Lemma~\ref{lem:commonmesh}, and the stage limit is the
least $J$ with $\frac78Lns^24^{-J}\le\varepsilon$ instead of
$\frac9{16}Lns^24^{-J}\le\varepsilon$, so a trial may run one stage longer.
\item Each trial starts with the center at the current incumbent instead
of~$\ell$. Validity and the rate analysis allow any first center (the bound
on $\norm{c-x^*}$ in (G2) of Lemma~\ref{lem:commonmesh} holds at $j=0$ for
every $c\in X$); the bit-length analysis of Appendix~\ref{app:growth} does
not cover it.
\item Before stage~$0$, two sweeps of exact coordinate descent from $\ell$,
$u$ and the box midpoint provide the first incumbent; after each stage the
same descent, within the current box, starts from the grid minimizer $y_j$.
An optional presolve certifies convex instances by an exact $LDL^\T$
factorization and a rational KKT point.
\item The lower bound used in the success test and reported at the end is the
largest bound of all stages and trials, including the initial interval bound,
rather than $\beta_j$ of the last stage.
\item The ``uniform'' grids of the experiments are the grids of
Definition~\ref{def:graded} with $\theta=0$ around one center, not the unions
of cells of UC (Algorithm~\ref{alg:uc}).
\item The exact-output wrapper, which implements EX (Algorithm~\ref{alg:ex}),
uses the acceptance rule of Proposition~\ref{prop:accept} with the row-sum
height constant of Remark~\ref{rem:heights}, with a common denominator that
also clears $H_{ij}/2$; this constant is never smaller than the one
in~\eqref{eq:exact-constants}. It starts at $q=4$, doubles $q$ after each
round, restarts the grid solver from the full box in every round with the
incumbent as its first center and, in place of $\ell$, as a start of the
descent, and tests three candidates after each round: the incumbent, its
coordinatewise rational reconstruction (Remark~\ref{rem:cf}) and the
stationary point of REC (Algorithm~\ref{alg:rec}).
\end{enumerate}
Further modules, each with its own checker, implement exact linear and convex
box-QP oracles with primal, dual or Farkas witnesses, affine recourse
recognition and lifting, minimum-cut and convex value-factor recourse, the
TU-constrained grids, the endpoint optimal-set descriptor and explicit
polynomial factors. The convex value-factor and recourse checkers reuse the
solver's grid, table and lifting code, so only the main checker is
independent in the sense of Section~\ref{sec:comp-impl}.

\paragraph{Planted family (Section~\ref{sec:comp-growth}).}
Lemma~\ref{lem:growthcert}(a) gives the certified lower bound $g_0$ on the
growth constant, by an exact $LDL^\T$ test of $H+2M-2g_0I$, and part~(b) with
a rational approximate eigenvector of $H_{J_0J_0}$ gives an upper bound; with
$L=2$ this brackets $\kappa$ to within $12\%$.

The following observation shows when a free coordinate of $x^*$ can be a
grid node. Within a trial every node of coordinate $i$ is either a dyadic rational or the
sum of the first center $c_i^{(0)}$ and a dyadic rational, because the
endpoints $\pm1$, the meshes and the gradings are dyadic, filtering moves
endpoints only to nodes, and every later center is a node. A coordinate
$x^*_i=k/21$ with $k\ne0$ is therefore a node only if $c_i^{(0)}$, which the
coordinate descent provides, differs from $x^*_i$ by a dyadic rational. In the
78 single-trial runs of E1 with filtering and of E3 with
$\kappa_{\mathrm{target}}=4$, 153 free coordinates of $x^*$ were nodes of the
last grid: 146 had already been set to $x^*_i$ by the initial descent, 3 had
$x^*_i=0$, and in the other 4 the first center differed from $x^*_i$ by a
nonzero dyadic rational.

\paragraph{Experiments E1--E3 (Section~\ref{sec:comp-growth}).}
The graded runs of E1 with filtering and the runs of E2 and E3 with the
theorem's grading satisfy $8\kappa\theta^2\le1$. At all 876 stages of these
runs, the gap $D(y_j)$ was at most $0.16\,Lnh_j^2$ (bound
$\frac9{16}Lnh_j^2$), the retained radius at most $0.18$ of
$4.2\sqrt{n\kappa}\,h_j$, and the number of nodes per coordinate at most
$5\%$ of the cap $8\theta^{-1}\lceil\log_2(n+2)\rceil$ (and at most $0.4\%$ of
the implementation's cap). The radius ratio uses the lower end of the
certified interval for $\kappa$, which can only increase it.
On the separable instances of E3 ($\kappa_{\mathrm{target}}=2$), filtered
uniform grids used exactly $4(\lfloor\sqrt{n/4}\rfloor+1)+1$ nodes;
Section~\ref{sec:comp-growth} reports the coupled case because the separable
one does not test coupling.

\paragraph{Exact output (Section~\ref{sec:comp-exact}).}
In E4, the 15 certified instances that needed more than one stage needed
67--77, 132--144, 275 and 534 stages for 49--65, 80--120, 208 and 342 bits,
with last-round precisions $q=64$, $128$, $256$ and $512$.

\paragraph{Comparison with SCIP (Section~\ref{sec:comp-scip}).}
Table~\ref{tab:scip} gives the results of E5 per instance class. The planted
objectives have constants between 94 and $1.6\cdot10^3$ and optimal value~$0$,
so the absolute target $10^{-6}$ is between $6\cdot10^{-10}$ and
$1.1\cdot10^{-8}$ of the constant; the inward derivatives at the active
coordinates are between 23 and 192.5. A rerun of the 21 SCIP runs reproduced
every primal bound. For SCIP's solution $(x,t)$, let $\bar x$ be the
projection of $x$ onto the box; then $F(\bar x)-t$ is the shortfall up to at
most $1.3\cdot10^{-14}$. The rerun evaluated $F$ exactly and split
$F(\bar x)-t$ into the bound part $F(\bar x)-F(x)$ and the epigraph part
$F(x)-t$. In every run the violated bounds were exactly those of the
coordinates at a bound in $x^*$, each violated outward by about $10^{-8}$, and
SCIP's own check accepted the solution. The bound part was $51\%$ to $94\%$ of
the shortfall, and the epigraph part was between $9.00\cdot10^{-7}$ and
$9.03\cdot10^{-7}$. In separate runs, giving SCIP 60
seconds (paths with $n\ge32$), setting its primal and dual feasibility
tolerances to $10^{-9}$ (paths with $n\ge32$), or moving the objective
constant into an offset (all 21 instances) did not change the outcome: SCIP
stopped at the time limit in all 27 runs on paths with $n\ge32$ (with 60
seconds, at dual bounds between $-7.5\cdot10^{-4}$ and $-5.1\cdot10^{-5}$);
with the offset, the exact gap on the 12 instances with $n\le16$ was still
$1.9\cdot10^{-6}$ to $5.8\cdot10^{-6}$; and its reported interval contained
the optimum in none of the 39 runs.

\begin{table}[t]
\centering\small
\setlength{\tabcolsep}{4pt}
\caption{Experiment E5: planted family, $\kappa\in[3.99,4.45]$, three seeds per
row. ``Bag'' is the maximum bag size. CT: maximum solve and replay times
(gap $\le10^{-6}$ certified on all instances, 60-second limit). SCIP: status,
maximum time (20-second limit), largest exact gap (exact value of SCIP's
incumbent, projected onto the box, minus its dual bound) and smallest dual
bound. Because $\OPT=0$ and SCIP's projected incumbents are optimal to within
$1.3\cdot10^{-14}$, the exact gap equals minus the dual bound to the digits
shown.}\label{tab:scip}
\begin{tabular}{lrrrrlrrr}
\toprule
&&&\multicolumn{2}{c}{CT (s)}&\multicolumn{4}{c}{SCIP}\\
\cmidrule(lr){4-5}\cmidrule(lr){6-9}
graph & $n$ & bag & solve & replay & status & time (s) & gap (exact) & dual bound\\
\midrule
path & 16 & 2 & 0.33 & 0.39 & gap reported & 5.4 & $2.9\cdot10^{-6}$ & $-2.9\cdot10^{-6}$\\
tree & 16 & 2 & 0.30 & 0.43 & gap reported & 0.8 & $5.5\cdot10^{-6}$ & $-5.5\cdot10^{-6}$\\
band & 12 & 3 & 1.33 & 2.12 & gap reported & 1.4 & $6.8\cdot10^{-6}$ & $-6.8\cdot10^{-6}$\\
band & 8 & 4 & 6.66 & 11.1 & gap reported & 0.3 & $5.5\cdot10^{-6}$ & $-5.5\cdot10^{-6}$\\
path & 32 & 2 & 0.99 & 1.20 & time limit & 20 & $1.4\cdot10^{-4}$ & $-1.4\cdot10^{-4}$\\
path & 64 & 2 & 3.02 & 3.62 & time limit & 20 & $3.0\cdot10^{-4}$ & $-3.0\cdot10^{-4}$\\
path & 128 & 2 & 6.90 & 9.39 & time limit & 20 & $7.6\cdot10^{-4}$ & $-7.6\cdot10^{-4}$\\
\bottomrule
\end{tabular}
\end{table}

\paragraph{Random instances (Section~\ref{sec:comp-random}).}
The instances of E6 have diagonal entries $k/2$ with $k$ uniform in
$\{-4,\dots,8\}$, interaction entries $\pm k/4$ with $k\in\{1,\dots,6\}$,
linear entries $k/4$ with $k\in\{-8,\dots,8\}$, and band width~2 for the band
instances. The candidate $\hat x$ is the stationary point of the face of the
final incumbent at $\varepsilon=2^{-50}$. When the sufficient condition of
Lemma~\ref{lem:growthcert}(a) holds at $\hat x$, it proves that $\hat x$ is
the unique minimizer and gives a lower bound $g_0$ on the point-growth
constant; an upper bound comes from part~(b) and from the definition of growth
at the points obtained by moving one active coordinate of $\hat x$ to its
other endpoint. With $L=\max_i\max\{H_{ii},0\}$ this brackets $\kappa$. The
condition held on 11 instances, with the certified intervals $[6.2,6.6]$,
$[11.6,15.7]$, $[14.0,22.3]$, $[6.1,24.4]$, $[10.0,28.6]$, $[8.1,31.6]$,
$[5.1,55.8]$, $[9.5,67.3]$, $[6.8,141.9]$, $[4.5,182.8]$ and
$[8.0,6662.2]$ for $\kappa$ (rounded outward). On the other 29 instances the
matrix $H+2M$ was not positive definite. On 28 of them the localized test of
Proposition~\ref{prop:local}, applied to the same certificate with the
stationary point of the face of the stage incumbent as candidate, proved that
candidate optimal; on one instance neither test succeeded. At
$\varepsilon=2^{-50}$ the 11 instances with certified growth had 9 to 11
nodes per coordinate, the other 29 had 8 to 18. The candidates are mostly off
the grids: at $\varepsilon=2^{-50}$, 42 of the 213 free coordinates of the
candidates were nodes of the last grid. Table~\ref{tab:random} gives the grid
sizes per instance class.

\begin{table}[t]
\centering\small
\caption{Experiment E6: CT on unplanted random box QPs, five instances per
row. ``Bag'' is the maximum bag size. ``Growth'' counts the instances on which
Lemma~\ref{lem:growthcert}(a) certifies point growth. ``Nodes'' is the largest
number of nodes of a coordinate grid over all stages and the five instances,
for each accuracy $\varepsilon$; ``stages'' and the time (maximum, seconds)
are for $\varepsilon=2^{-50}$.}\label{tab:random}
\begin{tabular}{lrrrrrrrrrr}
\toprule
&&&&\multicolumn{5}{c}{nodes at $\varepsilon=$}&&\\
\cmidrule(lr){5-9}
sparsity & $n$ & bag & growth & $2^{-10}$ & $2^{-20}$ & $2^{-30}$ & $2^{-40}$ & $2^{-50}$ & stages & time\\
\midrule
path & 8 & 2 & 3 & 16 & 16 & 16 & 16 & 16 & 27--28 & 0.20\\
path & 12 & 2 & 3 & 15 & 15 & 16 & 16 & 16 & 28 & 0.29\\
path & 16 & 2 & 0 & 12 & 13 & 15 & 15 & 15 & 28 & 0.42\\
path & 24 & 2 & 0 & 16 & 17 & 17 & 17 & 17 & 29 & 0.67\\
band & 8 & 3 & 2 & 10 & 10 & 10 & 10 & 10 & 27--28 & 0.57\\
band & 12 & 3 & 2 & 13 & 14 & 14 & 14 & 14 & 28 & 0.81\\
band & 16 & 3 & 1 & 15 & 16 & 18 & 18 & 18 & 28--29 & 5.22\\
band & 24 & 3 & 0 & 15 & 15 & 15 & 15 & 15 & 29 & 1.43\\
\bottomrule
\end{tabular}
\end{table}

\paragraph{Recourse (Section~\ref{sec:comp-recourse}).}
The affine stars have a core coordinate whose own term $-x_0^2$ is concave and
4 or 16 leaves coupled to it by stiff convex squares $64(y_i-x_0/2)^2$, with
the variables in random order; recognition finds the globally affine
responses, and the reduced problem is solved exactly. The piecewise instances
are the functions $\phi_M-z^2$ with $\phi_M$ from Example~\ref{ex:cv-fm}, for
$M=1$ and $M=100$. The dense instances ($n=7$ and $33$) have one convex core
coordinate and a residual that is concave in each coordinate and submodular,
with a complete interaction graph, so the decomposition used by plain grids
has a bag of size $n$. Table~\ref{tab:recourse} gives the results.

\begin{table}[t]
\centering\small
\setlength{\tabcolsep}{4pt}
\caption{Plain corrected grids (CT) versus the recourse pipeline (gap target
$2^{-10}$, 20-second limit). Status: ``exact'', exact minimizer certified;
``certified'', gap at most $2^{-10}$ certified; ``time limit'' and ``table
limit'', run stopped by the time limit or because the next stage would need
more than $2\cdot10^5$ table entries. ``Gap'' is the certified gap at the end
of the plain run, ``bag'' the maximum bag size of the decomposition used by
the plain grids, and ``entries'' the total number of table entries of the
plain run; times in seconds.}\label{tab:recourse}
\begin{tabular}{lrrlrrrrlrr}
\toprule
&&&\multicolumn{5}{c}{plain grids}&\multicolumn{3}{c}{recourse}\\
\cmidrule(lr){4-8}\cmidrule(lr){9-11}
instance & $n$ & bag & status & gap & entries & solve & replay & status & solve & replay\\
\midrule
affine star & 5 & 2 & time limit & 0.019 & 559{,}320 & 20.0 & 26.4 & exact & 0.005 & 0.001\\
affine star & 17 & 2 & time limit & 1.23 & 823{,}569 & 20.0 & 50.9 & exact & 0.061 & 0.057\\
piecewise, $M=1$ & 3 & 3 & certified & $4\cdot10^{-4}$ & 498 & 0.017 & 0.017 & certified & 0.013 & 0.003\\
piecewise, $M=100$ & 3 & 3 & certified & $7\cdot10^{-4}$ & 35{,}600 & 0.63 & 1.12 & certified & 0.004 & 0.002\\
dense cut & 7 & 7 & certified & $2^{-10}$ & 2{,}038 & 0.044 & 0.125 & exact & 0.003 & 0.003\\
dense cut & 33 & 33 & table limit & 2.96 & 0 & 0.021 & 0.009 & certified & 0.078 & 0.064\\
\bottomrule
\end{tabular}
\end{table}
